% Part: first-order-logic
% Chapter: completeness
% Section: introduction

\documentclass[../../../include/open-logic-section]{subfiles}

\begin{document}

\iftag{FOL}
      {\olfileid{fol}{com}{int}}
      {\olfileid{pl}{com}{int}}

\olsection{مقدمه}

قضیهٔ تمامیت یکی از بنیادی‌ترین نتایج منطق است. این قضیه دو صورت‌بندی
دارد که هم‌ارزی آن‌ها را اثبات خواهیم کرد. صورت‌بندی نخست مطلبی بنیادی
دربارهٔ رابطهٔ میان استلزام معنایی و دستگاه !!{derivation} ما می‌گوید:
اگر !!a{sentence}~$!A$ از چند !!{sentence} در~$\Gamma$ نتیجه شود، آنگاه
!!a{derivation} نیز وجود دارد که $\Gamma \Proves !A$ را اثبات می‌کند.
پس دستگاه !!{derivation} تا آنجا که ممکن است نیرومند است، بی‌آنکه
چیزهایی را اثبات کند که واقعاً نتیجه نمی‌شوند.

صورت‌بندی دوم را می‌توان نتیجه‌ای دربارهٔ وجود مدل دانست: هر مجموعهٔ
سازگار از !!{sentence}s ارضاپذیر است. سازگاری مفهومی برهان‌نظری است:
می‌گوید دستگاه !!{derivation} ما نمی‌تواند برخی !!{derivation}s را
پدید آورد. اما چه کسی می‌گوید صرفاً از اینکه هیچ !!{derivation} از نوعی
معین از~$\Gamma$ وجود ندارد، تضمین می‌شود که
\iftag{FOL}{!!a{structure}~$\Struct{M}$}{!!{valuation}~$\pAssign{v}$}
وجود دارد که $\iftag{FOL}{\Sat{M}}{\pSat{v}}{\Gamma}$؟ پیش از آنکه
قضیهٔ تمامیت نخستین‌بار اثبات شود---و در واقع پیش از آنکه دستگاه‌های
!!{derivation} امروزی را داشته باشیم---دیوید هیلبرت، ریاضی‌دان بزرگ
آلمانی، بر این باور بود که سازگاری نظریه‌های ریاضی وجود اشیایی را که
آن نظریه‌ها درباره‌شان سخن می‌گویند تضمین می‌کند. او در نامه‌ای به
گوتلوب فرگه چنین نوشت:
\begin{quote}
  اگر اصول موضوع دل‌بخواهیِ مفروض، با همهٔ نتایجشان، یکدیگر را نقض
  نکنند، آنگاه صادق‌اند و چیزهایی که اصول موضوع تعریف می‌کنند وجود
  دارند. این برای من معیار صدق و وجود است.
\end{quote}
فرگه به‌شدت مخالفت کرد. صورت‌بندی دوم قضیهٔ تمامیت نشان می‌دهد که هیلبرت
دست‌کم به این معنا حق داشت که اگر اصول موضوع سازگار باشند، آنگاه
\emph{یک}
\iftag{FOL}{!!{structure}}{!!{valuation}} وجود دارد که همهٔ آن‌ها را
صادق می‌سازد.

این‌ها تنها دلایل اهمیت قضیهٔ تمامیت---یا دقیق‌تر، برهان آن---نیستند.
این قضیه پیامدهای مهمی دارد که برخی از آن‌ها را جداگانه بررسی خواهیم
کرد. برای نمونه، چون هر !!{derivation} که $\Gamma \Proves !A$ را نشان
می‌دهد متناهی است و ازاین‌رو فقط می‌تواند شمار متناهی‌ای از
!!{sentence}s موجود در~$\Gamma$ را به کار ببرد، از قضیهٔ تمامیت نتیجه
می‌شود که اگر $!A$ پیامد~$\Gamma$ باشد، از پیش پیامد زیرمجموعه‌ای متناهی
از~$\Gamma$ است. این ویژگی را \emph{فشردگی} می‌نامند. به‌طور هم‌ارز، اگر
هر زیرمجموعهٔ متناهی~$\Gamma$ سازگار باشد، خود~$\Gamma$ نیز باید سازگار
باشد.

هرچند قضیهٔ فشردگی از راه میان‌بُرِ !!{derivation}s از قضیهٔ تمامیت
نتیجه می‌شود، می‌توان \emph{خود برهان} قضیهٔ تمامیت را نیز برای اثبات
مستقیم آن به کار برد. برهان، مجموعه‌ای از !!{sentence}s را که خاصیتی
معین---یعنی سازگاری---دارد می‌گیرد و از آن !!a{structure} با خواصی معین
می‌سازد (در اینجا، ساختاری که مجموعه را ارضا می‌کند). تقریباً همین ساخت
را می‌توان برای اثبات مستقیم فشردگی به کار برد؛ کافی است به‌جای
مجموعه‌های سازگار از !!{sentence}s، از مجموعه‌های «متناهی‌وار ارضاپذیر»
آغاز کنیم.
\iftag{FOL}{این ساخت پیامدهای دیگری نیز به دست می‌دهد؛ برای نمونه، هر
  مجموعهٔ ارضاپذیر از !!{sentence}s مدلی متناهی یا !!{denumerable} دارد.
  (این نتیجه قضیهٔ لوونهایم--اسکولم نامیده می‌شود.) به‌طور کلی، ساختن
  !!{structure}s از مجموعه‌های !!{sentence}s در منطق بسیار به کار
  می‌رود و گاه حتی در فلسفه نیز کاربرد دارد.}{}

\end{document}
